%% LyX 2.4.4 created this file.  For more info, see https://www.lyx.org/.
%% Do not edit unless you really know what you are doing.
\documentclass[english]{article}
\usepackage[T1]{fontenc}
\usepackage[latin9]{inputenc}
\usepackage{units}
\usepackage{enumitem}
\usepackage{amsmath}
\usepackage{amsthm}
\usepackage{amssymb}
\usepackage{geometry}
\geometry{verbose,tmargin=1in,bmargin=1in,lmargin=1in,rmargin=1in}

\makeatletter
%%%%%%%%%%%%%%%%%%%%%%%%%%%%%% Textclass specific LaTeX commands.
\numberwithin{equation}{section}
\numberwithin{figure}{section}
\newlength{\lyxlabelwidth}      % auxiliary length 

%%%%%%%%%%%%%%%%%%%%%%%%%%%%%% User specified LaTeX commands.
\usepackage{fancyhdr}
\usepackage{lastpage}
\pagestyle{fancy}
\usepackage{enumitem}
\usepackage{ifthen}
\everymath=\expandafter{\the\everymath\displaystyle}
%\DeclareMathSizes{10}{10}{10}{10}
\usepackage{multicol}

\usepackage{tikz}
\usepackage{tikz-3dplot}
\usetikzlibrary{calc,arrows}

\usetikzlibrary{patterns}
\usetikzlibrary{plotmarks}
\usepackage{pgfplots}
\pgfplotsset{compat=newest}

\usepackage{calc}
\usepackage[nomessages]{fp}

\usepackage{datetime}
% Added by lyx2lyx
\setlength{\parskip}{\smallskipamount}
\setlength{\parindent}{0pt}

\makeatother

\usepackage{babel}
\begin{document}
\renewcommand{\author}{Dr.\,James Ashe}

\newcommand{\classNum}{336}

\newcommand{\classSec}{A}

\newcommand{\examType}{HW 2.3}

\renewcommand{\date}{\formatdate{16}{10}{2013}}

%%First page's header%%%%%%%%%%%%%%%%%%%%%%
\fancypagestyle{firststyle} %%header settings for first pagr
{
	\fancyhead[L]{MTH \classNum{}(\classSec{}) \author{}\\
	\textbf{\examType{}} ~\date{}}
	\fancyhead[C]{}
	\fancyhead[R]{\textbf{Solutions}}
	\setlength{\headsep}{14pt}
}
%%%%%%%%%%%%%%%%%%%%%%%%%%%%%%%%%%%%%%%%%%

%%Next page's header%%%%%%%%%%%%%%%%%%%%%%%
\setlist{nolistsep}
\lhead{\classNum{}(\classSec{})} 
\chead{\textbf{Solutions}} 
\cfoot{\thepage{}~of~\pageref{LastPage}} 
\setlength{\headsep}{5pt} 
%%%%%%%%%%%%%%%%%%%%%%%%%%%%%%%%%%%%%%%%%%%

%%Various settings; see the preamble for other settings%%%%%
%\setenumerate[0]{leftmargin=12pt,itemindent=0pt,labelwidth=0pt}
\setlist{topsep=.5pt}
\renewcommand{\labelenumi}{\textbf{\arabic{enumi}.}}
\renewcommand{\labelenumii}{\textbf{\alph{enumii}.}}
\thispagestyle{firststyle}
%%%%%%%%%%%%%%%%%%%%%%%%%%%%%%%%%%%%%%%%%%

\newcommand{\xmax}{0}
\newcommand{\xmin}{-\xmax}
\newcommand{\xstep}{0}
\newcommand{\ymax}{0}
\newcommand{\ymin}{-\ymax}
\newcommand{\ystep}{0} 
\newcommand{\dspace}{\vspace{1cm}}
\newcommand{\xa}{0}
\newcommand{\xb}{0}
\newcommand{\ya}{1}
\newcommand{\yb}{1}

\small
\begin{enumerate}[resume]
\item \vspace{0cm}

\begin{enumerate}[resume]
\item [\textbf{b.}]$\begin{bmatrix}\begin{array}{rr|rr}
1 & 3 & 1 & 0\\
2 & 6 & 0 & 1
\end{array}\end{bmatrix}\rightsquigarrow\begin{array}{r}
\\\mbox{R2-2R1}
\end{array}\begin{bmatrix}\begin{array}{rr|rr}
1 & 3 & 1 & 0\\
0 & 0 & -2 & 1
\end{array}\end{bmatrix}$ so $A^{-1}$ does not exist.
\end{enumerate}
\item \vspace{0cm}

\begin{enumerate}[resume]
\item \vspace{0cm}

\begin{enumerate}
\item $\begin{bmatrix}\begin{array}{rr|rr}
2 & 3 & 1 & 0\\
3 & 5 & 0 & 1
\end{array}\end{bmatrix}\rightsquigarrow\begin{array}{r}
\\\mbox{2R2}
\end{array}\begin{bmatrix}\begin{array}{rr|rr}
2 & 3 & 1 & 0\\
6 & 10 & 0 & 2
\end{array}\end{bmatrix}\rightsquigarrow\begin{array}{r}
\\\mbox{R2-3R1}
\end{array}\begin{bmatrix}\begin{array}{rr|rr}
2 & 3 & 1 & 0\\
0 & 1 & -3 & 2
\end{array}\end{bmatrix}\rightsquigarrow\begin{array}{r}
\mbox{R1-3R2}\\
\mbox{}
\end{array}\begin{bmatrix}\begin{array}{rr|rr}
2 & 0 & 10 & -6\\
0 & 1 & -3 & 2
\end{array}\end{bmatrix}\rightsquigarrow\begin{array}{r}
\nicefrac{1}{2}\mbox{R1}\\
\mbox{}
\end{array}\begin{bmatrix}\begin{array}{rr|rr}
1 & 0 & 5 & -3\\
0 & 1 & -3 & 2
\end{array}\end{bmatrix}$ So $A^{-1}=\begin{bmatrix}\begin{array}{rr}
5 & -3\\
-3 & 2
\end{array}\end{bmatrix}$
\item 
\begin{eqnarray*}
A^{-1}\left(A\mathbf{x}\right) & = & A^{-1}\mathbf{b}\\
\mathbf{x} & = & \begin{bmatrix}\begin{array}{rr}
5 & -3\\
-3 & 2
\end{array}\end{bmatrix}\begin{bmatrix}\begin{array}{r}
3\\
4
\end{array}\end{bmatrix}\\
\mathbf{x} & = & \begin{bmatrix}\begin{array}{r}
3\\
-1
\end{array}\end{bmatrix}
\end{eqnarray*}
\item $\mathbf{b=}3\begin{bmatrix}\begin{array}{r}
2\\
3
\end{array}\end{bmatrix}-1\begin{bmatrix}\begin{array}{r}
3\\
5
\end{array}\end{bmatrix}=\begin{bmatrix}\begin{array}{r}
3\\
4
\end{array}\end{bmatrix}$.
\end{enumerate}
\item [\textbf{c.}]

\begin{enumerate}
\item $\begin{bmatrix}\begin{array}{rrr|rrr}
1 & 1 & 1 & 1 & 0 & 0\\
0 & 1 & 1 & 0 & 1 & 0\\
1 & 2 & 1 & 0 & 0 & 1
\end{array}\end{bmatrix}\rightsquigarrow\begin{array}{r}
\mbox{R1-R2}\\
\\\mbox{R3-R1}
\end{array}\begin{bmatrix}\begin{array}{rrr|rrr}
1 & 0 & 0 & 1 & -1 & 0\\
0 & 1 & 1 & 0 & 1 & 0\\
0 & 1 & 0 & -1 & 0 & 1
\end{array}\end{bmatrix}\rightsquigarrow\begin{array}{r}
\\\mbox{R2-R3}\\
\\\end{array}\begin{bmatrix}\begin{array}{rrr|rrr}
1 & 0 & 0 & 1 & -1 & 0\\
0 & 0 & 1 & 1 & 1 & -1\\
0 & 1 & 0 & -1 & 0 & 1
\end{array}\end{bmatrix}\rightsquigarrow\begin{array}{r}
\\\\\mbox{}
\end{array}\begin{bmatrix}\begin{array}{rrr|rrr}
1 & 0 & 0 & 1 & -1 & 0\\
0 & 1 & 0 & -1 & 0 & 1\\
0 & 0 & 1 & 1 & 1 & -1
\end{array}\end{bmatrix}$
\[
A^{-1}=\begin{bmatrix}\begin{array}{rrr}
1 & -1 & 0\\
-1 & 0 & 1\\
1 & 1 & -1
\end{array}\end{bmatrix}
\]
\end{enumerate}
\begin{enumerate}[resume]
\item 
\begin{eqnarray*}
A^{-1}\left(A\mathbf{x}\right) & = & A^{-1}\mathbf{b}\\
\mathbf{x} & = & \begin{bmatrix}\begin{array}{rrr}
1 & -1 & 0\\
-1 & 0 & 1\\
1 & 1 & -1
\end{array}\end{bmatrix}\begin{bmatrix}\begin{array}{r}
3\\
0\\
1
\end{array}\end{bmatrix}=\begin{bmatrix}\begin{array}{r}
3\\
-2\\
2
\end{array}\end{bmatrix}
\end{eqnarray*}
\item 
\[
\mathbf{b}=3\begin{bmatrix}\begin{array}{r}
1\\
0\\
1
\end{array}\end{bmatrix}-2\begin{bmatrix}\begin{array}{r}
1\\
1\\
2
\end{array}\end{bmatrix}+2\begin{bmatrix}\begin{array}{r}
1\\
1\\
1
\end{array}\end{bmatrix}
\]
\end{enumerate}
\end{enumerate}
\item \vspace{0cm}

\begin{enumerate}
\item $\left(BAB^{-1}\right)^{2}=BAB^{-1}BAB^{-1}=BAAB^{-1}=BA^{2}B^{-1}$
\end{enumerate}
\begin{enumerate}[resume]
\item {[}not assigned{]} $\left(BAB^{-1}\right)^{n}=BA\left(B^{-1}B\right)A\left(B^{-1}B\right)A\left(B^{-1}B\right)AB^{-1}\cdots BA\left(B^{-1}B\right)AB^{-1}=BA^{n}B^{-1}$
\item By proposition 3.4 $(BAB^{-1})^{-1}=\left(B^{-1}\right)^{-1}\left(BA\right)^{-1}=B\left(A^{-1}B^{-1}\right)=BA^{-1}B^{-1}$
but we must assume that $A$ and $B$ are invertible.\\
Note that this means part (b) holds for negative integers as well.
\end{enumerate}
\item $A$ is invertible so there exists an $A^{-1}$. So then 
\begin{eqnarray*}
A\mathbf{x} & = & 7\mathbf{x}\\
A^{-1}A\mathbf{x} & = & A^{-1}7\mathbf{x}\\
\mathbf{x} & = & A^{-1}7\mathbf{x}\\
\frac{1}{7}\mathbf{x} & = & A^{-1}\mathbf{x}
\end{eqnarray*}
$\Rightarrow A^{-1}\mathbf{x}=\frac{1}{7}\mathbf{x}$.
\end{enumerate}

\end{document}
